\documentclass[10pt,a4paper]{amsart}
\usepackage[margin=19mm]{geometry}
\usepackage{amsmath,amssymb,mathtools}
\usepackage[scr=rsfs,cal=euler]{mathalfa}
\usepackage{xfrac}
\usepackage[unicode,hidelinks,bookmarks=false]{hyperref}
\newcommand{\ie}{i.e.\ }
\newcommand{\cf}{cf.\ }
\newcommand{\resp}{respectively\ }
\newcommand{\Subsection}{\subsection}
\providecommand{\nobreakdash}{\nobreak\hbox{-}}
\setlength{\emergencystretch}{2em}
\multlinegap0pt
\usepackage[shortlabels]{enumitem}
\usepackage{cancel}
\usepackage{amssymb}
\usepackage{xcolor}
\newtheorem{lemma}{Lemma}
\newcommand{\tr}{\triangleright}
\title{Generalized digroups, di-skew braces, and solutions of the set-theoretic Yang--Baxter equation}
\author{Andrea ALBANO, Paola STEFANELLI}
\date{Source: arXiv:2505.15387v2 (CC BY 4.0)}
\begin{document}
\maketitle

\noindent\emph{Source and licence.} Generalized digroups, di-skew braces, and solutions of the set-theoretic Yang--Baxter equation. Version arXiv:2505.15387v2 (CC BY 4.0), distributed by arXiv under the Creative Commons Attribution 4.0 International licence, http://creativecommons.org/licenses/by/4.0/, as recorded in the arXiv licence metadata for this exact version and shown on the abstract page https://arxiv.org/abs/2505.15387v2. Full licence notice: \texttt{licenses/arxiv-2505.15387-CC-BY-4.0.txt}.\par\smallskip

\noindent\textit{Editorial scope.} Counted unit: the article's lemma powers\_prop (source0.tex lines 537--544). The article's complete original proof of it is delivered with it (source0.tex lines 545--574, one \texttt{proof} environment of 30 lines). No mathematics is written, repaired or re-derived: the statement and the proof are the article's own lines, in the article's own notation, and no retained line is substituted or re-flowed.  Retained uncounted as the source-local context the proof needs: lines 429--431.  Nothing was omitted from any retained span, so the package carries no omission record.  No retained line contains a citation, so the wrapper carries no bibliography and no bibliography entry.  No retained line contains a figure environment, an image-inclusion command or drawing code, so no image is delivered and the package declares no asset dependency.  Editorial formatting only, in addition: an AMS \texttt{amsart} preamble that restates the article's own declarations and macros used by the retained lines, namely the environments \texttt{lemma} on separate counters, together with the standard packages those lines need.  The retained environments print in this document as Lemma 1 in order of appearance, whereas the article prints the same items with the numbering of its own full document.  The complete native source of the article as distributed in the arXiv e-print for this exact version is bundled unchanged as \texttt{sources/178507/source0.tex}.
\begin{align}\label{self_distri1}
    \forall\, a,b \in D \quad L_aL_b = L_{L_a(b)}L_a \,,
\end{align}

\begin{lemma}\label{powers_prop}
    Let $(D,\triangleright)$ be a rack. Then, the following identities 
    \begin{align}
        r_\tr^{2n}(a,b) &= \left((L_bL_a)^nL_a^{-n}(a),(L_bL_a)^nL_b^{-n}(b)\right) \label{even_powers2}\,, \\[0.2cm]
        r_\tr^{2n+1}(a,b) &= \left((L_bL_a)^nL_b^{-n}(b),(L_bL_a)^{n+1}L_a^{-(n+1)}(a)\right) \label{odd_powers2}\,,
    \end{align}
    hold, for all $a,b \in D$ and $n\in\mathbb{N}$.
\end{lemma}

\begin{proof}
    As a technical identity, let us first note that, {by \eqref{self_distri1}},
    \begin{align}\label{eq:tech_id}
        L_{\left(L_xL_y\right)^kL_u^{-k}(v)} =
        \left(L_xL_y\right)^kL_u^{-k}L_vL_u^k\left(L_xL_y\right)^{-k}
    \end{align}
    holds for all $x,y,u,v \in D$ and $k \in \mathbb{N}$.
    To prove \eqref{even_powers2} we proceed by induction on $n \in \mathbb{N}$.
    In the base case $n = 1$ we have that 
    $r^2_\tr(a,b) = \left(L_b(a),L_{L_b(a)}(b)\right) = \left(L_b(a),L_bL_aL_b^{-1}(b)\right)$, for all $a,b \in D$.
    Now, let $n \in \mathbb{N}$ with $n > 1$ and assume the assertion true for $n > 1$.
    Fix $a,b \in D$ and note that
    \begin{align*}
        r_\tr^{2(n+1)}(a,b) 
            = 
            r_\tr^2 r_\tr^{2n}(a,b)
            =
            (L_B(A),L_{L_B(A)}(B)) \,,
    \end{align*}
    where we have set $A = (L_bL_a)^{n}L_a^{-n}(a)$ and $B = (L_bL_a)^{n}L_b^{-n}(b)$.
    Applying \eqref{eq:tech_id}, we compute $L_B(A) = \left(L_bL_a\right)^{n+1}L_a^{-n-1}(a)$ and again by \eqref{eq:tech_id} deduce that
    \begin{align*}
        L_{L_B(A)}(B) 
        = \left(L_bL_a\right)^{n+1}L_a\left(L_bL_a\right)^{-n-1}(B) 
        % &= \left(L_bL_a\right)^{n+1}L_a\left(L_bL_a\right)^{-1}L_b^{-n}(b) \\
        = \left(L_bL_a\right)^{n+1}L_b^{-n-1}(b) \,,
    \end{align*}
    thus proving the assertion.
    To verify that \eqref{odd_powers2} holds, for all $n \in \mathbb{N}$, it is sufficient to compute $r_\tr^{2n+1} = r^{}_\tr r_\tr^{2n}$ using \eqref{even_powers2} and \eqref{eq:tech_id} accordingly.
\end{proof}

\end{document}
